\section{Post-hoc analyses}
\label{sec:s-posthoc}
The analyses in this section were specified after the primary results were known, committed (\texttt{configs/\allowbreak posthoc.json}) before they were run, and use the splits, inner folds, selection objective, checkpoint and threshold rules, cross-fitting recipe and final seeds of the primary protocol. No primary result was changed. Sections~S5.7--S5.9 were added after a pre-submission review, when all other results were known (\texttt{configs/\allowbreak posthoc.json}, key \texttt{review\_round\_2}); they use only the saved per-sample predictions and selected configurations.

\subsection{Exposure-time baselines: selected configurations}
Time-LR searched all six values of $C$; Time+game HGB used the HGB search space and budget (16 configurations).

\begin{center}\footnotesize
\begin{tabular}{lcllc}
\toprule
Model & Outer fold & Feature set & Selected hyper-parameters & Inner-OOF cost \\
\midrule
Time-LR & 0 & time & C = 0.001 & 0.471 \\
Time-LR & 1 & time & C = 1 & 0.528 \\
Time-LR & 2 & time & C = 0.1 & 0.549 \\
Time-LR & 3 & time & C = 0.01 & 0.497 \\
Time-LR & 4 & time & C = 0.1 & 0.550 \\
Time+game HGB & 0 & time+game & lr 0.02, iter 200, leaves 8, min leaf 300, $\ell_2$ 1 & 0.480 \\
Time+game HGB & 1 & time+game & lr 0.02, iter 200, leaves 8, min leaf 300, $\ell_2$ 0 & 0.501 \\
Time+game HGB & 2 & time+game & lr 0.02, iter 100, leaves 4, min leaf 300, $\ell_2$ 10 & 0.524 \\
Time+game HGB & 3 & time+game & lr 0.02, iter 200, leaves 4, min leaf 300, $\ell_2$ 1 & 0.522 \\
Time+game HGB & 4 & time+game & lr 0.02, iter 200, leaves 8, min leaf 300, $\ell_2$ 0 & 0.506 \\
\bottomrule
\end{tabular}

\end{center}

\subsection{Holm adjustment of the paired comparisons}
Group-bootstrap comparisons of E-DDQN with the six other primary models under each operating rule, with Holm's step-down adjustment over the six comparisons of each rule.

\begin{center}\footnotesize
\begin{tabular}{llcccc}
\toprule
Rule & B & $\Delta$Cost (E-DDQN$-$B) & 95\,\% CI & $p_{\text{group}}$ & $p_{\text{Holm}}$ \\
\midrule
D1 (min cost) & E-CB & +0.175 & [+0.016, +0.448] & 0.014 & 0.070 \\
D1 (min cost) & DQN & +0.190 & [$-$0.017, +0.493] & 0.079 & 0.105 \\
D1 (min cost) & CS-MLP & +0.190 & [+0.005, +0.485] & 0.039 & 0.105 \\
D1 (min cost) & LogReg & +0.194 & [+0.011, +0.477] & 0.026 & 0.104 \\
D1 (min cost) & HGB & +0.183 & [+0.011, +0.460] & 0.035 & 0.105 \\
D1 (min cost) & RF & +0.242 & [+0.052, +0.528] & 0.001 & 0.006 \\
D0 ($\hat p\ge1/6$) & E-CB & +0.214 & [+0.025, +0.487] & 0.018 & 0.054 \\
D0 ($\hat p\ge1/6$) & DQN & +0.242 & [+0.039, +0.542] & 0.012 & 0.048 \\
D0 ($\hat p\ge1/6$) & CS-MLP & +0.204 & [$-$0.030, +0.507] & 0.113 & 0.226 \\
D0 ($\hat p\ge1/6$) & LogReg & +0.087 & [$-$0.293, +0.541] & 0.694 & 0.694 \\
D0 ($\hat p\ge1/6$) & HGB & +0.328 & [+0.078, +0.639] & 0.003 & 0.015 \\
D0 ($\hat p\ge1/6$) & RF & +0.411 & [+0.147, +0.748] & 0.001 & 0.006 \\
val.\ recall $\ge$0.95 & E-CB & +0.220 & [+0.017, +0.566] & 0.021 & 0.080 \\
val.\ recall $\ge$0.95 & DQN & +0.236 & [+0.013, +0.580] & 0.028 & 0.080 \\
val.\ recall $\ge$0.95 & CS-MLP & +0.294 & [+0.033, +0.652] & 0.013 & 0.065 \\
val.\ recall $\ge$0.95 & LogReg & +0.307 & [+0.027, +0.687] & 0.020 & 0.080 \\
val.\ recall $\ge$0.95 & HGB & +0.314 & [+0.030, +0.711] & 0.023 & 0.080 \\
val.\ recall $\ge$0.95 & RF & +0.342 & [+0.067, +0.719] & 0.003 & 0.018 \\
\bottomrule
\end{tabular}

\end{center}

\subsection{Comparison with always alarming}
Paired bootstrap of each model's D1 cost (seed ensemble) minus the cost of always alarming; Holm adjustment over the nine models (group units).

\begin{center}\footnotesize
\resizebox{\textwidth}{!}{\begin{tabular}{lcccccc}
\toprule
Model & $\Delta$Cost (model$-$always) & CI (groups) & $p$ & $p_{\text{Holm}}$ & CI (sessions) & $p$ \\
\midrule
E-DDQN & +0.168 & [$-$0.024, +0.466] & 0.104 & 0.728 & [$-$0.026, +0.438] & 0.116 \\
E-CB & $-$0.007 & [$-$0.093, +0.103] & 0.810 & 1.000 & [$-$0.089, +0.105] & 0.800 \\
DQN & $-$0.022 & [$-$0.104, +0.062] & 0.560 & 1.000 & [$-$0.099, +0.075] & 0.598 \\
CS-MLP & $-$0.022 & [$-$0.046, $-$0.002] & 0.037 & 0.296 & [$-$0.049, +0.004] & 0.109 \\
LogReg & $-$0.026 & [$-$0.061, +0.006] & 0.136 & 0.816 & [$-$0.058, +0.005] & 0.113 \\
HGB & $-$0.015 & [$-$0.122, +0.121] & 0.765 & 1.000 & [$-$0.104, +0.092] & 0.744 \\
RF & $-$0.074 & [$-$0.135, $-$0.019] & 0.012 & 0.108 & [$-$0.120, $-$0.025] & 0.002 \\
Time-LR & $-$0.010 & [$-$0.073, +0.080] & 0.737 & 1.000 & [$-$0.080, +0.080] & 0.782 \\
Time+game HGB & $-$0.041 & [$-$0.100, +0.051] & 0.302 & 1.000 & [$-$0.099, +0.048] & 0.275 \\
\bottomrule
\end{tabular}
}
\end{center}

\subsection{Grouping sensitivity}
Paired bootstrap resampling sessions instead of linkage groups (Holm over the six primary comparisons), and the range of the cost difference in a leave-one-group-out jackknife over the \nGroupsEval{} groups.

\begin{center}\footnotesize
\resizebox{\textwidth}{!}{\begin{tabular}{lcccccc}
\toprule
B & $\Delta$Cost (E-DDQN$-$B) & CI (sessions) & $p$ & $p_{\text{Holm}}$ & jackknife range & sign flips \\
\midrule
E-CB & +0.175 & [+0.014, +0.423] & 0.014 & 0.070 & +0.068 to +0.201 & 0 \\
DQN & +0.190 & [+0.001, +0.454] & 0.050 & 0.087 & +0.084 to +0.248 & 0 \\
CS-MLP & +0.190 & [+0.006, +0.452] & 0.037 & 0.087 & +0.084 to +0.225 & 0 \\
LogReg & +0.194 & [+0.007, +0.454] & 0.029 & 0.087 & +0.094 to +0.225 & 0 \\
HGB & +0.183 & [+0.014, +0.440] & 0.021 & 0.084 & +0.076 to +0.215 & 0 \\
RF & +0.242 & [+0.055, +0.503] & 0.001 & 0.006 & +0.137 to +0.284 & 0 \\
Time-LR & +0.178 & [$-$0.015, +0.439] & 0.089 & -- & +0.075 to +0.213 & 0 \\
Time+game HGB & +0.209 & [+0.016, +0.470] & 0.030 & -- & +0.101 to +0.248 & 0 \\
\bottomrule
\end{tabular}
}
\end{center}

\subsection{Validation-chosen cost curves}
Test cost of the seed ensemble divided by the cost of the better trivial policy, for each false-negative to false-positive ratio $R$ (threshold chosen on inner-validation predictions at each $R$). At $R=5$ the values reproduce the D1 seed-ensemble costs exactly (checked programmatically).

\begin{center}\footnotesize
\resizebox{\textwidth}{!}{\begin{tabular}{lcccccccccc}
\toprule
Model & 1 & 1.5 & 2 & 3 & 4 & 5 & 7 & 10 & 15 & 20 \\
\midrule
E-DDQN & 1.004 & 0.921 & 1.034 & 1.130 & 1.241 & 1.301 & 1.382 & 1.322 & 1.448 & 1.508 \\
E-CB & 0.828 & 0.819 & 0.900 & 0.920 & 0.954 & 0.988 & 1.068 & 0.992 & 1.007 & 1.013 \\
DQN & 0.921 & 0.796 & 0.865 & 0.845 & 0.929 & 0.961 & 1.013 & 1.068 & 1.208 & 1.304 \\
CS-MLP & 0.846 & 0.773 & 0.835 & 0.980 & 0.960 & 0.961 & 0.969 & 0.981 & 0.985 & 0.985 \\
LogReg & 0.852 & 0.868 & 0.915 & 0.941 & 0.946 & 0.953 & 0.966 & 0.983 & 0.989 & 1.010 \\
HGB & 0.897 & 0.770 & 0.808 & 0.884 & 0.945 & 0.973 & 1.001 & 0.980 & 0.991 & 0.998 \\
RF & 0.861 & 0.717 & 0.773 & 0.836 & 0.864 & 0.868 & 0.893 & 0.928 & 0.960 & 0.983 \\
Time-LR & 0.843 & 0.718 & 0.781 & 0.840 & 0.883 & 0.982 & 0.990 & 0.984 & 1.000 & 1.000 \\
Time+game HGB & 0.816 & 0.731 & 0.760 & 0.833 & 0.882 & 0.926 & 0.942 & 1.014 & 1.053 & 1.064 \\
\bottomrule
\end{tabular}
}
\end{center}

\subsection{Q-value scale}
Largest absolute Q-value over the final runs of each outer fold relative to the return bound $r_{\max}/(1-\gamma)$ of the selected discount ($r_{\max}=0.5$ for the asymmetric reward, 1 for the symmetric reward).

\begin{center}\footnotesize
\begin{tabular}{lccccc}
\toprule
Model & Outer fold & $\gamma$ & max $|Q|$ & bound $r_{\max}/(1-\gamma)$ & ratio \\
\midrule
E-DDQN & 0 & 0.99 & 41.45 & 50.00 & 0.83 \\
E-DDQN & 1 & 0.99 & 61.78 & 50.00 & 1.24 \\
E-DDQN & 2 & 0.99 & 35.08 & 50.00 & 0.70 \\
E-DDQN & 3 & 0.99 & 4.13 & 50.00 & 0.08 \\
E-DDQN & 4 & 0.9 & 11.56 & 5.00 & 2.31 \\
E-CB & 0 & 0 & 0.70 & 0.50 & 1.41 \\
E-CB & 1 & 0 & 0.68 & 0.50 & 1.37 \\
E-CB & 2 & 0 & 0.68 & 0.50 & 1.36 \\
E-CB & 3 & 0 & 0.89 & 0.50 & 1.79 \\
E-CB & 4 & 0 & 0.95 & 0.50 & 1.89 \\
DQN & 0 & 0 & 2.24 & 1.00 & 2.24 \\
DQN & 1 & 0.9 & 4.09 & 10.00 & 0.41 \\
DQN & 2 & 0.9 & 3.05 & 10.00 & 0.30 \\
DQN & 3 & 0.9 & 10.10 & 10.00 & 1.01 \\
DQN & 4 & 0 & 2.08 & 1.00 & 2.08 \\
\bottomrule
\end{tabular}

\end{center}

\subsection{Within-fold ROC-AUC}
Pooled ROC-AUC combines the test scores of five fold-specific models and therefore mixes discrimination with differences in score scale between folds. Within-fold ROC-AUC is computed in each outer-test fold and averaged over folds (seed mean: also averaged over the five seeds).

\begin{center}\footnotesize
\begin{tabular}{lcccc}
\toprule
Model & Pooled (seed mean) & Within-fold (seed mean) & Pooled (ensemble) & Within-fold (ensemble) \\
\midrule
E-DDQN & 0.538 & 0.539 & 0.541 & 0.552 \\
E-CB & 0.653 & 0.658 & 0.659 & 0.665 \\
DQN & 0.661 & 0.666 & 0.665 & 0.666 \\
CS-MLP & 0.644 & 0.666 & 0.654 & 0.672 \\
LogReg & 0.626 & 0.636 & 0.626 & 0.636 \\
HGB & 0.692 & 0.702 & 0.692 & 0.702 \\
RF & 0.688 & 0.715 & 0.688 & 0.715 \\
Time-LR & 0.704 & 0.709 & 0.704 & 0.709 \\
Time+game HGB & 0.676 & 0.699 & 0.676 & 0.699 \\
\bottomrule
\end{tabular}

\end{center}

\subsection{Exact group sign-flip tests}
Two-sided sign-flip permutation test of the pooled D1 cost difference of the seed ensembles, with the \nGroupsEval{} linkage groups as units and all $2^{21}$ sign patterns enumerated. Holm adjustment over the six primary comparisons with E-DDQN and over the nine comparisons with always alarming. The group-bootstrap $p$-values of Section~S5.2 (rule D1) are shown for comparison.

\begin{center}\footnotesize
\resizebox{\textwidth}{!}{\begin{tabular}{lccccc}
\toprule
Comparison (D1, seed ensembles) & $\Delta$Cost & $p_{\text{boot}}$ & $p_{\text{boot,Holm}}$ & $p_{\text{sign-flip}}$ & $p_{\text{sign-flip,Holm}}$ \\
\midrule
E-DDQN vs E-CB & +0.175 & 0.014 & 0.070 & 0.048 & 0.241 \\
E-DDQN vs DQN & +0.190 & 0.079 & 0.105 & 0.127 & 0.330 \\
E-DDQN vs CS-MLP & +0.190 & 0.039 & 0.105 & 0.116 & 0.330 \\
E-DDQN vs LogReg & +0.194 & 0.026 & 0.104 & 0.094 & 0.330 \\
E-DDQN vs HGB & +0.183 & 0.035 & 0.105 & 0.083 & 0.330 \\
E-DDQN vs RF & +0.242 & 0.001 & 0.006 & 0.014 & 0.083 \\
E-DDQN vs Time-LR & +0.178 & -- & -- & 0.117 & -- \\
E-DDQN vs Time+game HGB & +0.209 & -- & -- & 0.067 & -- \\
E-DDQN vs always alarm & +0.168 &  &  & 0.250 & 1.000 \\
E-CB vs always alarm & $-$0.007 &  &  & 0.922 & 1.000 \\
DQN vs always alarm & $-$0.022 &  &  & 0.620 & 1.000 \\
CS-MLP vs always alarm & $-$0.022 &  &  & 0.053 & 0.422 \\
LogReg vs always alarm & $-$0.026 &  &  & 0.172 & 1.000 \\
HGB vs always alarm & $-$0.015 &  &  & 0.769 & 1.000 \\
RF vs always alarm & $-$0.074 &  &  & 0.022 & 0.194 \\
Time-LR vs always alarm & $-$0.010 &  &  & 0.820 & 1.000 \\
Time+game HGB vs always alarm & $-$0.041 &  &  & 0.358 & 1.000 \\
\bottomrule
\end{tabular}
}
\end{center}

\subsection{Selected feature sets and within-fold ROC-AUC}
Exploratory. Feature set selected on validation data for each primary model and outer fold (T: telemetry; T+Q: telemetry and questionnaire indicators; E: exposure; E+Q: exposure and questionnaire indicators) and the within-fold ROC-AUC of that model on the fold's test groups (mean over seeds). Feature sets were selected, not assigned, so the association does not establish a cause.

\begin{center}\footnotesize
\begin{tabular}{lccccc}
\toprule
Model & Fold 0 & Fold 1 & Fold 2 & Fold 3 & Fold 4 \\
\midrule
E-DDQN & E+Q 0.541 & E+Q 0.505 & E+Q 0.723 & E+Q 0.450 & E+Q 0.476 \\
E-CB & E 0.642 & E 0.717 & E 0.748 & E 0.737 & E+Q 0.446 \\
DQN & E 0.706 & E+Q 0.530 & E 0.843 & E 0.736 & E+Q 0.515 \\
CS-MLP & E+Q 0.562 & E+Q 0.607 & E 0.812 & T 0.702 & T 0.644 \\
LogReg & E 0.623 & T+Q 0.428 & E 0.810 & E 0.712 & E+Q 0.610 \\
HGB & E 0.631 & E 0.740 & E 0.791 & T 0.761 & E 0.585 \\
RF & T 0.630 & E 0.728 & E 0.771 & E 0.761 & T+Q 0.683 \\
\bottomrule
\end{tabular}

\end{center}

